\documentclass[11pt,a4paper]{article}
\usepackage[utf8]{inputenc}
\usepackage[T1]{fontenc}
\usepackage[brazil]{babel}
\usepackage[margin=2.5cm]{geometry}
\usepackage{mathptmx}
\usepackage{enumitem}
\usepackage{booktabs}
\usepackage{tabularx}
\usepackage{array}
\usepackage{hyperref}
\hypersetup{colorlinks=false, pdfborder={0 0 0}}
\usepackage{fancyhdr}

\pagestyle{fancy}
\fancyhf{}
\fancyhead[L]{\small TrackFleet --- Métricas e Sinal de Alerta}
\fancyhead[R]{\small\thepage}
\renewcommand{\headrulewidth}{0.4pt}

\newcolumntype{L}{>{\raggedright\arraybackslash}X}

\begin{document}

\begin{center}
{\Large\bfseries DOCUMENTO DE GOVERNANÇA DO PROJETO}\\[2pt]
{\large\bfseries TrackFleet --- Gestão de Rastreador de Automóveis}\\[4pt]
Disciplina de Gerenciamento de Projetos $\cdot$ Framework PMBOK\textregistered\ 8\textsuperscript{a} Edição $\cdot$ Domínio: Governança\\[2pt]
Integrantes: Enzo Allebrand e Kerlon Ribeiro (dois GPs) $\cdot$ 4\textsuperscript{o} semestre $\cdot$ Data: 13/08
\end{center}

\vspace{2mm}
\hrule
\vspace{4mm}

\noindent\textbf{\large Painel de Métricas e Sinal de Alerta}

\vspace{1mm}
\noindent Como o projeto vai saber que está entregando valor, e qual sinal antecipa um problema antes de ele virar crise.

\section{Métricas de Resultado}
Medem valor entregue, não esforço.
\begin{itemize}[itemsep=3pt]
  \item \textbf{Resultado 1:} número de veículos ativamente monitorados por semana.
  \item \textbf{Resultado 2:} percentual de alertas de manutenção atendidos antes de virarem problema, ou seja, manutenção preventiva em vez de corretiva.
  \item \textbf{Resultado 3:} redução percentual estimada no custo de manutenção corretiva, reportada pelas transportadoras piloto.
\end{itemize}

\section{Sinal de Alerta}
Indicador antecedente, que avisa cedo se algo vai mal:
\begin{itemize}[itemsep=3pt]
  \item Menos de três veículos cadastrados e ativos até a segunda semana de testes aciona o alerta de baixa adoção.
\end{itemize}

\section{O que Não Usar como Métrica}
É esforço, não resultado:
\begin{itemize}[itemsep=3pt]
  \item Número de linhas de código escritas.
\end{itemize}

\end{document}
